% Options for packages loaded elsewhere
% Options for packages loaded elsewhere
\PassOptionsToPackage{unicode}{hyperref}
\PassOptionsToPackage{hyphens}{url}
%
\documentclass[
  11pt,
  ignorenonframetext,
  aspectratio=169,
]{beamer}
\newif\ifbibliography
\usepackage{pgfpages}
\setbeamertemplate{caption}[numbered]
\setbeamertemplate{caption label separator}{: }
\setbeamercolor{caption name}{fg=normal text.fg}
\beamertemplatenavigationsymbolsempty
% remove section numbering
\setbeamertemplate{part page}{
  \centering
  \begin{beamercolorbox}[sep=16pt,center]{part title}
    \usebeamerfont{part title}\insertpart\par
  \end{beamercolorbox}
}
\setbeamertemplate{section page}{
  \centering
  \begin{beamercolorbox}[sep=12pt,center]{section title}
    \usebeamerfont{section title}\insertsection\par
  \end{beamercolorbox}
}
\setbeamertemplate{subsection page}{
  \centering
  \begin{beamercolorbox}[sep=8pt,center]{subsection title}
    \usebeamerfont{subsection title}\insertsubsection\par
  \end{beamercolorbox}
}
% Prevent slide breaks in the middle of a paragraph
\widowpenalties 1 10000
\raggedbottom
\AtBeginPart{
  \frame{\partpage}
}
\AtBeginSection{
  \ifbibliography
  \else
    \frame{\sectionpage}
  \fi
}
\AtBeginSubsection{
  \frame{\subsectionpage}
}
\usepackage{iftex}
\ifPDFTeX
  \usepackage[T1]{fontenc}
  \usepackage[utf8]{inputenc}
  \usepackage{textcomp} % provide euro and other symbols
\else % if luatex or xetex
  \usepackage{unicode-math} % this also loads fontspec
  \defaultfontfeatures{Scale=MatchLowercase}
  \defaultfontfeatures[\rmfamily]{Ligatures=TeX,Scale=1}
\fi
\usepackage{lmodern}

\ifPDFTeX\else
  % xetex/luatex font selection
\fi
% Use upquote if available, for straight quotes in verbatim environments
\IfFileExists{upquote.sty}{\usepackage{upquote}}{}
\IfFileExists{microtype.sty}{% use microtype if available
  \usepackage[]{microtype}
  \UseMicrotypeSet[protrusion]{basicmath} % disable protrusion for tt fonts
}{}
\makeatletter
\@ifundefined{KOMAClassName}{% if non-KOMA class
  \IfFileExists{parskip.sty}{%
    \usepackage{parskip}
  }{% else
    \setlength{\parindent}{0pt}
    \setlength{\parskip}{6pt plus 2pt minus 1pt}}
}{% if KOMA class
  \KOMAoptions{parskip=half}}
\makeatother

\usepackage{color}
\usepackage{fancyvrb}
\newcommand{\VerbBar}{|}
\newcommand{\VERB}{\Verb[commandchars=\\\{\}]}
\DefineVerbatimEnvironment{Highlighting}{Verbatim}{commandchars=\\\{\}}
% Add ',fontsize=\small' for more characters per line
\usepackage{framed}
\definecolor{shadecolor}{RGB}{241,243,245}
\newenvironment{Shaded}{\begin{snugshade}}{\end{snugshade}}
\newcommand{\AlertTok}[1]{\textcolor[rgb]{0.68,0.00,0.00}{#1}}
\newcommand{\AnnotationTok}[1]{\textcolor[rgb]{0.37,0.37,0.37}{#1}}
\newcommand{\AttributeTok}[1]{\textcolor[rgb]{0.40,0.45,0.13}{#1}}
\newcommand{\BaseNTok}[1]{\textcolor[rgb]{0.68,0.00,0.00}{#1}}
\newcommand{\BuiltInTok}[1]{\textcolor[rgb]{0.00,0.23,0.31}{#1}}
\newcommand{\CharTok}[1]{\textcolor[rgb]{0.13,0.47,0.30}{#1}}
\newcommand{\CommentTok}[1]{\textcolor[rgb]{0.37,0.37,0.37}{#1}}
\newcommand{\CommentVarTok}[1]{\textcolor[rgb]{0.37,0.37,0.37}{\textit{#1}}}
\newcommand{\ConstantTok}[1]{\textcolor[rgb]{0.56,0.35,0.01}{#1}}
\newcommand{\ControlFlowTok}[1]{\textcolor[rgb]{0.00,0.23,0.31}{\textbf{#1}}}
\newcommand{\DataTypeTok}[1]{\textcolor[rgb]{0.68,0.00,0.00}{#1}}
\newcommand{\DecValTok}[1]{\textcolor[rgb]{0.68,0.00,0.00}{#1}}
\newcommand{\DocumentationTok}[1]{\textcolor[rgb]{0.37,0.37,0.37}{\textit{#1}}}
\newcommand{\ErrorTok}[1]{\textcolor[rgb]{0.68,0.00,0.00}{#1}}
\newcommand{\ExtensionTok}[1]{\textcolor[rgb]{0.00,0.23,0.31}{#1}}
\newcommand{\FloatTok}[1]{\textcolor[rgb]{0.68,0.00,0.00}{#1}}
\newcommand{\FunctionTok}[1]{\textcolor[rgb]{0.28,0.35,0.67}{#1}}
\newcommand{\ImportTok}[1]{\textcolor[rgb]{0.00,0.46,0.62}{#1}}
\newcommand{\InformationTok}[1]{\textcolor[rgb]{0.37,0.37,0.37}{#1}}
\newcommand{\KeywordTok}[1]{\textcolor[rgb]{0.00,0.23,0.31}{\textbf{#1}}}
\newcommand{\NormalTok}[1]{\textcolor[rgb]{0.00,0.23,0.31}{#1}}
\newcommand{\OperatorTok}[1]{\textcolor[rgb]{0.37,0.37,0.37}{#1}}
\newcommand{\OtherTok}[1]{\textcolor[rgb]{0.00,0.23,0.31}{#1}}
\newcommand{\PreprocessorTok}[1]{\textcolor[rgb]{0.68,0.00,0.00}{#1}}
\newcommand{\RegionMarkerTok}[1]{\textcolor[rgb]{0.00,0.23,0.31}{#1}}
\newcommand{\SpecialCharTok}[1]{\textcolor[rgb]{0.37,0.37,0.37}{#1}}
\newcommand{\SpecialStringTok}[1]{\textcolor[rgb]{0.13,0.47,0.30}{#1}}
\newcommand{\StringTok}[1]{\textcolor[rgb]{0.13,0.47,0.30}{#1}}
\newcommand{\VariableTok}[1]{\textcolor[rgb]{0.07,0.07,0.07}{#1}}
\newcommand{\VerbatimStringTok}[1]{\textcolor[rgb]{0.13,0.47,0.30}{#1}}
\newcommand{\WarningTok}[1]{\textcolor[rgb]{0.37,0.37,0.37}{\textit{#1}}}

\usepackage{longtable,booktabs,array}
\newcounter{none} % for unnumbered tables
\usepackage{calc} % for calculating minipage widths
\usepackage{caption}
% Make caption package work with longtable
\makeatletter
\def\fnum@table{\tablename~\thetable}
\makeatother
\usepackage{graphicx}
\makeatletter
\newsavebox\pandoc@box
\newcommand*\pandocbounded[1]{% scales image to fit in text height/width
  \sbox\pandoc@box{#1}%
  \Gscale@div\@tempa{\textheight}{\dimexpr\ht\pandoc@box+\dp\pandoc@box\relax}%
  \Gscale@div\@tempb{\linewidth}{\wd\pandoc@box}%
  \ifdim\@tempb\p@<\@tempa\p@\let\@tempa\@tempb\fi% select the smaller of both
  \ifdim\@tempa\p@<\p@\scalebox{\@tempa}{\usebox\pandoc@box}%
  \else\usebox{\pandoc@box}%
  \fi%
}
% Set default figure placement to htbp
\def\fps@figure{htbp}
\makeatother





\setlength{\emergencystretch}{3em} % prevent overfull lines

\providecommand{\tightlist}{%
  \setlength{\itemsep}{0pt}\setlength{\parskip}{0pt}}



 


\setbeamertemplate{navigation symbols}{}
\setbeamertemplate{footline}[frame number]
\setbeamertemplate{itemize items}[circle]
\setbeamertemplate{itemize subitem}[triangle]
\setbeamercolor{frametitle}{fg=black}
\setbeamerfont{frametitle}{series=\bfseries}
\newcommand{\indep}{\perp\!\!\!\perp}
\usepackage{booktabs}
% shrink code input (fancyvrb) and code output (verbatim)
\usepackage{fancyvrb}
\fvset{fontsize=\scriptsize}
\makeatletter
\def\verbatim@font{\scriptsize\ttfamily}
\makeatother
\makeatletter
\@ifpackageloaded{caption}{}{\usepackage{caption}}
\AtBeginDocument{%
\ifdefined\contentsname
  \renewcommand*\contentsname{Table of contents}
\else
  \newcommand\contentsname{Table of contents}
\fi
\ifdefined\listfigurename
  \renewcommand*\listfigurename{List of Figures}
\else
  \newcommand\listfigurename{List of Figures}
\fi
\ifdefined\listtablename
  \renewcommand*\listtablename{List of Tables}
\else
  \newcommand\listtablename{List of Tables}
\fi
\ifdefined\figurename
  \renewcommand*\figurename{Figure}
\else
  \newcommand\figurename{Figure}
\fi
\ifdefined\tablename
  \renewcommand*\tablename{Table}
\else
  \newcommand\tablename{Table}
\fi
}
\@ifpackageloaded{float}{}{\usepackage{float}}
\floatstyle{ruled}
\@ifundefined{c@chapter}{\newfloat{codelisting}{h}{lop}}{\newfloat{codelisting}{h}{lop}[chapter]}
\floatname{codelisting}{Listing}
\newcommand*\listoflistings{\listof{codelisting}{List of Listings}}
\makeatother
\makeatletter
\makeatother
\makeatletter
\@ifpackageloaded{caption}{}{\usepackage{caption}}
\@ifpackageloaded{subcaption}{}{\usepackage{subcaption}}
\makeatother

\usepackage{bookmark}
\IfFileExists{xurl.sty}{\usepackage{xurl}}{} % add URL line breaks if available
\urlstyle{same}
\hypersetup{
  pdftitle={Causality and Potential Outcomes},
  pdfauthor={Professor Ji-Woong Chung},
  hidelinks,
  pdfcreator={LaTeX via pandoc}}


\title{\texorpdfstring{Causality and Potential
Outcomes}{Causality and Potential Outcomes}}
\subtitle{\texorpdfstring{BUSS975: Causal Inference in Financial
Research}{BUSS975: Causal Inference in Financial Research}}
\author{Professor Ji-Woong Chung}
\date{}
\institute{Korea University}

\begin{document}
\frame{\titlepage}


\begin{frame}[fragile]{Acknowledgments and readings}
\protect\phantomsection\label{acknowledgments-and-readings}
\begin{itemize}
\tightlist
\item
  These slides are based on Scott Cunningham, \emph{Causal Inference:
  The Remix} (2nd ed.), Chapter 4, ``Potential Outcomes and
  Randomization.''

  \begin{itemize}
  \tightlist
  \item
    Free online: \texttt{https://mixtape.scunning.com}
  \end{itemize}
\item
  The regression material in Part 6 is adapted from Todd Gormley's
  lecture notes.
\end{itemize}
\end{frame}

\begin{frame}{Outline}
\protect\phantomsection\label{outline}
\begin{enumerate}
\tightlist
\item
  What causality is --- and isn't
\item
  The potential outcomes framework
\item
  Selection bias: the Perfect Regulator
\item
  Randomization and independence
\item
  SUTVA
\item
  The regression bridge
\item
  Randomization inference
\end{enumerate}
\end{frame}

\section{What Causality Is --- and
Isn't}\label{what-causality-is-and-isnt}

\begin{frame}{Motivation}
\protect\phantomsection\label{motivation}
\begin{itemize}
\tightlist
\item
  As researchers, we want to make \textbf{causal} statements:

  \begin{itemize}
  \tightlist
  \item
    What is the effect of a corporate tax change on firms' leverage?
  \item
    What is the effect of CEO equity ownership on risk-taking?
  \item
    What is the effect of a disclosure mandate on the cost of capital?
  \end{itemize}
\item
  We are rarely satisfied with saying variables are merely
  ``associated.''
\item
  But nearly all of our data is \textbf{observational}: nobody
  randomized leverage, ownership, or regulation for us.
\item
  This course is about when --- and how --- observational data can
  answer causal questions.
\end{itemize}
\end{frame}

\begin{frame}{Three goals of empirical analysis}
\protect\phantomsection\label{three-goals-of-empirical-analysis}
\begin{enumerate}
\tightlist
\item
  \textbf{Description} --- how things are or were.

  \begin{itemize}
  \tightlist
  \item
    Has industry concentration increased over time?
  \end{itemize}
\item
  \textbf{Prediction} --- guessing an unknown value, without
  intervening.

  \begin{itemize}
  \tightlist
  \item
    What will this firm's earnings be next quarter?
  \end{itemize}
\item
  \textbf{Causality} --- how \emph{changing} one variable would change
  another, all else equal.

  \begin{itemize}
  \tightlist
  \item
    How would a disclosure mandate change firms' cost of capital?
  \end{itemize}
\end{enumerate}

\vspace{0.5em}

\begin{itemize}
\tightlist
\item
  All three are legitimate. But they require different tools, and
  confusing them is the original sin of empirical work.
\end{itemize}
\end{frame}

\begin{frame}{Fallacy 1: prediction is not causation}
\protect\phantomsection\label{fallacy-1-prediction-is-not-causation}
\begin{itemize}
\tightlist
\item
  Momentum: past returns \emph{predict} future returns.

  \begin{itemize}
  \tightlist
  \item
    Does buying past winners \emph{cause} them to keep winning? No ---
    the predictive relationship says nothing about intervention.
  \end{itemize}
\item
  A model can predict superbly while being causally empty.

  \begin{itemize}
  \tightlist
  \item
    This is the promise (and the limit) of much of machine learning.
  \end{itemize}
\item
  \textbf{Causal inference asks a different question}: what would happen
  if we \emph{changed} something?

  \begin{itemize}
  \tightlist
  \item
    Prediction asks what \(Y\) \textbf{is}, among the units where
    \(X = x\).
  \item
    Causation asks what \(Y\) \textbf{would be}, if we set \(X = x\).
  \end{itemize}
\end{itemize}
\end{frame}

\begin{frame}{Fallacy 2: temporal ordering is not causation}
\protect\phantomsection\label{fallacy-2-temporal-ordering-is-not-causation}
\begin{itemize}
\tightlist
\item
  \emph{Post hoc ergo propter hoc}: ``after this, therefore because of
  this.''
\item
  The rooster crows, then the sun rises.
\item
  Finance version: a struggling firm fires its CEO, then performance
  recovers.

  \begin{itemize}
  \tightlist
  \item
    Did the new CEO cause the recovery?
  \item
    Or is this \textbf{mean reversion} --- firms replace CEOs at the
    trough?
  \end{itemize}
\item
  Timing evidence is useful (we will use it in event studies and DiD),
  but timing alone proves nothing.
\end{itemize}
\end{frame}

\begin{frame}{Fallacy 3: no correlation does not mean no causation}
\protect\phantomsection\label{fallacy-3-no-correlation-does-not-mean-no-causation}
\begin{itemize}
\tightlist
\item
  A sailor moves the rudder constantly, yet the boat sails straight.

  \begin{itemize}
  \tightlist
  \item
    Correlation between rudder and course: \(\approx 0\). Causal effect:
    enormous.
  \item
    Why? The rudder \emph{responds optimally} to the wind.
  \end{itemize}
\item
  Finance is full of optimizing agents doing exactly this:

  \begin{itemize}
  \tightlist
  \item
    Firms hedge \textbf{because} they are exposed: hedging may show zero
    correlation with distress even if it strongly prevents distress.
  \item
    The Fed adjusts rates to offset shocks: policy can look uncorrelated
    with output.
  \end{itemize}
\item
  \textbf{Lesson}: when treatment is chosen optimally, causal effects
  can be \emph{invisible} in correlations.
\end{itemize}
\end{frame}

\begin{frame}{The common thread: the assignment mechanism}
\protect\phantomsection\label{the-common-thread-the-assignment-mechanism}
\begin{itemize}
\tightlist
\item
  In all three fallacies, what breaks the link between correlation and
  causation is \textbf{how treatment got assigned}:

  \begin{itemize}
  \tightlist
  \item
    Who chose to be treated? Why? Based on what?
  \end{itemize}
\item
  The central message of this lecture (and this course):
\end{itemize}

\begin{center}
\emph{Correlations are properties of the data.\\
Causal effects are properties of the data \textbf{plus} the assignment mechanism.}
\end{center}

\begin{itemize}
\tightlist
\item
  To formalize this, we need a language for ``what would have happened
  otherwise.''
\end{itemize}
\end{frame}

\begin{frame}{Defining causality: counterfactuals}
\protect\phantomsection\label{defining-causality-counterfactuals}
\begin{itemize}
\tightlist
\item
  Our (narrow) definition: \emph{ceteris paribus}, how does \(D\) affect
  \(Y\)?
\item
  To know the effect of \(D\) on unit \(i\), we need to compare:

  \begin{itemize}
  \tightlist
  \item
    \(Y_i\) \textbf{with} the treatment, and
  \item
    \(Y_i\) \textbf{without} the treatment,
  \item
    with everything else the same.
  \end{itemize}
\item
  One of these is always hypothetical --- a \textbf{counterfactual}.
\item
  \textbf{The Fundamental Problem of Causal Inference} (Holland 1986):
  we can never observe both states for the same unit at the same time.
\end{itemize}
\end{frame}

\section{The Potential Outcomes
Framework}\label{the-potential-outcomes-framework}

\begin{frame}{States of the world}
\protect\phantomsection\label{states-of-the-world}
\begin{itemize}
\tightlist
\item
  Imagine each unit exists in two potential states of the world,
  identical in every respect except one: treatment.
\item
  Binary treatment indicator: \[
  D_i = \begin{cases} 1 & \text{if unit } i \text{ is treated} \\ 0 & \text{otherwise} \end{cases}
  \]
\item
  Two \textbf{potential outcomes} for every unit:

  \begin{itemize}
  \tightlist
  \item
    \(Y_i^1\): the outcome unit \(i\) \emph{would} have if treated,
  \item
    \(Y_i^0\): the outcome unit \(i\) \emph{would} have if untreated.
  \end{itemize}
\item
  Both are defined for every unit --- treated or not. That is the trick.
\end{itemize}
\end{frame}

\begin{frame}{The switching equation}
\protect\phantomsection\label{the-switching-equation}
\begin{itemize}
\tightlist
\item
  What we \emph{observe} is only one of the two: \[
  Y_i = D_i\, Y_i^1 + (1 - D_i)\, Y_i^0
  \]
\item
  Treatment ``switches on'' which potential outcome is realized.
\item
  Notice what this equation does \textbf{not} contain: no error term, no
  functional form, no linearity. It is a definition, not a model.
\end{itemize}
\end{frame}

\begin{frame}{Individual treatment effects}
\protect\phantomsection\label{individual-treatment-effects}
\begin{itemize}
\tightlist
\item
  The causal effect of the treatment on unit \(i\): \[
  \delta_i = Y_i^1 - Y_i^0
  \]
\item
  \(\delta_i\) carries an \(i\) subscript.
\item
  Effects are allowed to \textbf{differ across units} from the start.
\item
  The Fundamental Problem, restated: \(\delta_i\) is never observed for
  any unit, because one of its two ingredients is always missing.

  \begin{itemize}
  \tightlist
  \item
    Causal inference is a \textbf{missing data problem} --- but the
    missing value never existed to begin with.
  \end{itemize}
\end{itemize}
\end{frame}

\begin{frame}{A running example: capital injections}
\protect\phantomsection\label{a-running-example-capital-injections}
\begin{itemize}
\tightlist
\item
  Setting: a banking crisis. The government can inject capital into
  banks (think TARP, 2008).
\item
  Outcome \(Y\): bank lending growth (\%) over the following year.
\item
  \(Y_i^1\): bank \(i\)'s lending growth \emph{with} an injection.
\item
  \(Y_i^0\): bank \(i\)'s lending growth \emph{without} one.
\item
  Question: what is the causal effect of capital injections on lending?
\end{itemize}
\end{frame}

\begin{frame}{Ten banks, two potential outcomes}
\protect\phantomsection\label{ten-banks-two-potential-outcomes}
{\def\LTcaptype{none} % do not increment counter
\begin{longtable}[]{@{}cccc@{}}
\toprule\noalign{}
Bank & \(Y^1\) & \(Y^0\) & \(\delta\) \\
\midrule\noalign{}
\endhead
1 & 9 & 2 & 7 \\
2 & 6 & 1 & 5 \\
3 & 7 & 3 & 4 \\
4 & 5 & 2 & 3 \\
5 & 6 & 4 & 2 \\
6 & 7 & 6 & 1 \\
7 & 7 & 7 & 0 \\
8 & 7 & 8 & \(-1\) \\
9 & 7 & 9 & \(-2\) \\
10 & 9 & 8 & 1 \\
\bottomrule\noalign{}
\end{longtable}
}

\begin{itemize}
\tightlist
\item
  Weak banks (low \(Y^0\)) gain a lot; healthy banks gain little --- or
  are even hurt.
\end{itemize}
\end{frame}

\begin{frame}{Average treatment effects}
\protect\phantomsection\label{average-treatment-effects}
\begin{itemize}
\tightlist
\item
  \textbf{Average Treatment Effect (ATE)}: \[
  ATE = E[\delta_i] = E[Y_i^1] - E[Y_i^0]
  \]
\item
  In our table: \(ATE = \frac{7+5+4+3+2+1+0-1-2+1}{10} = 2.0\)
\item
  On average, injections raise lending growth by 2 percentage points.
\item
  But note the heterogeneity: effects range from \(+7\) to \(-2\).

  \begin{itemize}
  \tightlist
  \item
    An average can hide winners and losers. Remember this in week 10
    (complex DiD).
  \end{itemize}
\end{itemize}
\end{frame}

\begin{frame}{ATT and ATU: different questions, different answers}
\protect\phantomsection\label{att-and-atu-different-questions-different-answers}
\begin{itemize}
\tightlist
\item
  \textbf{Average Treatment effect on the Treated}: \[
  ATT = E[\delta_i \mid D_i = 1]
  \]
\item
  \textbf{Average Treatment effect on the Untreated}: \[
  ATU = E[\delta_i \mid D_i = 0]
  \]
\item
  These are \textbf{different policy questions}:

  \begin{itemize}
  \tightlist
  \item
    ATT: ``Did TARP help the banks that received it?''
  \item
    ATE: ``Would injections help a randomly chosen bank?''
  \item
    ATU: ``Would extending the program to the remaining banks help
    them?''
  \end{itemize}
\item
  When effects are heterogeneous and assignment is not random,
  \(ATT \neq ATE \neq ATU\) --- and each estimator you will meet in this
  course targets a \emph{particular} one of these.
\end{itemize}
\end{frame}

\begin{frame}{What we actually observe}
\protect\phantomsection\label{what-we-actually-observe}
{\def\LTcaptype{none} % do not increment counter
\begin{longtable}[]{@{}ccccc@{}}
\toprule\noalign{}
Bank & \(D\) & \(Y^1\) & \(Y^0\) & \(\delta\) \\
\midrule\noalign{}
\endhead
1 & 1 & 9 & ? & ? \\
2 & 1 & 6 & ? & ? \\
3 & 1 & 7 & ? & ? \\
4 & 1 & 5 & ? & ? \\
5 & 1 & 6 & ? & ? \\
6 & 0 & ? & 6 & ? \\
7 & 0 & ? & 7 & ? \\
8 & 0 & ? & 8 & ? \\
9 & 0 & ? & 9 & ? \\
10 & 0 & ? & 8 & ? \\
\bottomrule\noalign{}
\end{longtable}
}

\begin{itemize}
\tightlist
\item
  Half of the table is missing --- forever. Everything we do from here
  on is a strategy for filling in the question marks \emph{on average}.
\end{itemize}
\end{frame}

\section{Selection Bias: The Perfect
Regulator}\label{selection-bias-the-perfect-regulator}

\begin{frame}{The naive comparison}
\protect\phantomsection\label{the-naive-comparison}
\begin{itemize}
\tightlist
\item
  With observational data, the natural move is to compare treated and
  untreated: \[
  SDO = E[Y_i^1 \mid D_i = 1] - E[Y_i^0 \mid D_i = 0]
  \]
\item
  \textbf{Simple Difference in Outcomes} (SDO): mean outcome of the
  treated minus mean outcome of the untreated.
\item
  Every cross-sectional regression of \(Y\) on \(D\) is, at heart, an
  SDO.
\item
  Question: when does SDO recover the ATE?
\end{itemize}
\end{frame}

\begin{frame}{The Perfect Regulator}
\protect\phantomsection\label{the-perfect-regulator}
\begin{itemize}
\tightlist
\item
  Suppose the regulator has perfect foresight: it knows every bank's
  \(\delta_i\) and injects capital into the five banks that \textbf{gain
  the most} (banks 1--5).
\item
  This is not a villain --- it is an \emph{optimizing} regulator,
  allocating scarce capital where it helps most.

  \begin{itemize}
  \tightlist
  \item
    Exactly what a good policymaker (or a good CEO, or a good doctor)
    should do.
  \end{itemize}
\item
  Treated: banks 1--5. Untreated: banks 6--10.
\item
  Now compute the SDO from the observable data.
\end{itemize}
\end{frame}

\begin{frame}{The Perfect Regulator's data}
\protect\phantomsection\label{the-perfect-regulators-data}
\begin{itemize}
\tightlist
\item
  Mean outcome among treated:
  \(E[Y^1 \mid D=1] = \frac{9+6+7+5+6}{5} = 6.6\)
\item
  Mean outcome among untreated:
  \(E[Y^0 \mid D=0] = \frac{6+7+8+9+8}{5} = 7.6\) \[
  SDO = 6.6 - 7.6 = -1.0
  \]
\item
  The truth: \(ATE = +2.0\), and \(ATT = \frac{7+5+4+3+2}{5} = +4.2\).
\item
  The naive comparison gets the \textbf{sign wrong}: injections look
  \emph{harmful} precisely because they were allocated \emph{well}.
\item
  This is the TARP debate in one table: ``bailed-out banks lent less
  than healthy banks'' is an SDO, not a causal effect.
\end{itemize}
\end{frame}

\begin{frame}{The decomposition}
\protect\phantomsection\label{the-decomposition}
\begin{itemize}
\item
  The SDO can always be decomposed --- this is an identity, not an
  assumption (\(\pi\) = share of units treated): \[
  \begin{aligned}
  \underbrace{E[Y^1 \mid D{=}1] - E[Y^0 \mid D{=}0]}_{\text{SDO}}
  = ATE &+ \underbrace{E[Y^0 \mid D{=}1] - E[Y^0 \mid D{=}0]}_{\text{selection bias}} \\
  &+ \underbrace{(1-\pi)\,(ATT - ATU)}_{\text{heterogeneous TE bias}}
  \end{aligned}
  \]
\item
  \textbf{Selection bias}: treated and untreated differ in their
  \emph{untreated} outcomes.
\item
  \textbf{Heterogeneous treatment effect bias}: the treated gain
  differently than the untreated would.
\item
  Memorize this equation. Every research design in this course is an
  attempt to kill these two terms.
\end{itemize}
\end{frame}

\begin{frame}{Proving the decomposition}
\protect\phantomsection\label{proving-the-decomposition}
Let \(\pi = P(D{=}1)\). Two steps, no assumptions.

\textbf{Step 1.} Add and subtract \(E[Y^0 \mid D{=}1]\) --- the
\emph{counterfactual} mean for the treated: \[
\begin{aligned}
SDO &= E[Y^1 \mid D{=}1] - E[Y^0 \mid D{=}0] \\
    &= \underbrace{E[Y^1 \mid D{=}1] - E[Y^0 \mid D{=}1]}_{ATT} \;+\;
       \underbrace{E[Y^0 \mid D{=}1] - E[Y^0 \mid D{=}0]}_{\text{selection bias}}
\end{aligned}
\]

\textbf{Step 2.} By the law of total expectation,
\(ATE = \pi\,ATT + (1-\pi)\,ATU\). Subtract \(ATE\) from \(ATT\): \[
ATT - ATE \;=\; (1-\pi)\,ATT - (1-\pi)\,ATU \;=\; (1-\pi)(ATT - ATU)
\]

Substituting \(ATT = ATE + (1-\pi)(ATT-ATU)\) into Step 1 gives the
three-term decomposition. \(\blacksquare\)
\end{frame}

\begin{frame}{The decomposition, in our example}
\protect\phantomsection\label{the-decomposition-in-our-example}
\begin{itemize}
\tightlist
\item
  Selection bias: \[
  E[Y^0 \mid D{=}1] - E[Y^0 \mid D{=}0] = \frac{2+1+3+2+4}{5} - \frac{6+7+8+9+8}{5} = 2.4 - 7.6 = -5.2
  \]

  \begin{itemize}
  \tightlist
  \item
    Treated banks were \emph{much weaker} to begin with.
  \end{itemize}
\item
  Heterogeneous TE bias, with \(\pi = 0.5\): \[
  (1 - 0.5)(ATT - ATU) = 0.5 \times (4.2 - (-0.2)) = +2.2
  \]
\item
  Check: \(\;2.0 - 5.2 + 2.2 = -1.0 = SDO\) \checkmark
\item
  Both bias terms are large; they partially offset, and the net result
  is a sign flip.
\end{itemize}
\end{frame}

\begin{frame}{Selection on gains is endemic in finance}
\protect\phantomsection\label{selection-on-gains-is-endemic-in-finance}
\begin{itemize}
\tightlist
\item
  The Perfect Doctor (Cunningham's version) treats patients who benefit
  most. Finance is full of Perfect Doctors:

  \begin{itemize}
  \tightlist
  \item
    Regulators bail out the banks where intervention matters most.
  \item
    Boards grant equity to the CEOs they expect to respond.
  \item
    Firms go public when going public helps them most.
  \item
    Banks lend to the borrowers they screen as creditworthy.
  \end{itemize}
\item
  Agents \textbf{optimize}. Optimization means treatment is assigned on
  expected gains --- the worst case for naive comparisons.
\item
  This is why finance cannot simply borrow ``big data'' logic: more
  observations do not fix a biased assignment mechanism.
\end{itemize}
\end{frame}

\begin{frame}[fragile]{Simulation: triage vs.~coin flip}
\protect\phantomsection\label{simulation-triage-vs.-coin-flip}
\begin{Shaded}
\begin{Highlighting}[]
\FunctionTok{set.seed}\NormalTok{(}\DecValTok{975}\NormalTok{)}
\NormalTok{n    }\OtherTok{\textless{}{-}} \DecValTok{100000}
\NormalTok{y0   }\OtherTok{\textless{}{-}} \FunctionTok{rnorm}\NormalTok{(n, }\DecValTok{7}\NormalTok{, }\DecValTok{2}\NormalTok{)     }\CommentTok{\# lending growth, no injection}
\NormalTok{gain }\OtherTok{\textless{}{-}} \FunctionTok{rnorm}\NormalTok{(n, }\DecValTok{2}\NormalTok{, }\DecValTok{1}\NormalTok{)     }\CommentTok{\# true individual gain (so ATE = 2)}
\NormalTok{y1   }\OtherTok{\textless{}{-}}\NormalTok{ y0 }\SpecialCharTok{+}\NormalTok{ gain}

\NormalTok{D1 }\OtherTok{\textless{}{-}} \FunctionTok{as.integer}\NormalTok{(y0 }\SpecialCharTok{\textless{}} \FunctionTok{quantile}\NormalTok{(y0, }\FloatTok{0.30}\NormalTok{))  }\CommentTok{\# triage: weakest 30\%}
\NormalTok{D2 }\OtherTok{\textless{}{-}} \FunctionTok{rbinom}\NormalTok{(n, }\DecValTok{1}\NormalTok{, }\FloatTok{0.30}\NormalTok{)                   }\CommentTok{\# coin flip: random 30\%}

\NormalTok{sdo }\OtherTok{\textless{}{-}} \ControlFlowTok{function}\NormalTok{(D) }\FunctionTok{mean}\NormalTok{(y1[D }\SpecialCharTok{==} \DecValTok{1}\NormalTok{]) }\SpecialCharTok{{-}} \FunctionTok{mean}\NormalTok{(y0[D }\SpecialCharTok{==} \DecValTok{0}\NormalTok{])}
\FunctionTok{round}\NormalTok{(}\FunctionTok{c}\NormalTok{(}\AttributeTok{ATE =} \FunctionTok{mean}\NormalTok{(gain), }\AttributeTok{triage =} \FunctionTok{sdo}\NormalTok{(D1), }\AttributeTok{coin =} \FunctionTok{sdo}\NormalTok{(D2)), }\DecValTok{2}\NormalTok{)}
\end{Highlighting}
\end{Shaded}

\begin{verbatim}
   ATE triage   coin 
  2.00  -1.31   1.98 
\end{verbatim}
\end{frame}

\begin{frame}{What the simulation tells us}
\protect\phantomsection\label{what-the-simulation-tells-us}
\begin{itemize}
\tightlist
\item
  The true effect is \(+2\) in both worlds --- same banks, same gains,
  only the \textbf{assignment rule} differs.
\item
  Triage (treat the weakest): SDO is strongly \emph{negative}. An
  observer concludes bailouts destroy lending.
\item
  Coin flip: SDO \(\approx\) ATE. Same data-generating process, correct
  answer.
\item
  With \(n = 100{,}000\): the bias does not shrink with sample size.
  \textbf{Selection bias is not a small-sample problem.}
\end{itemize}
\end{frame}

\section{Randomization and
Independence}\label{randomization-and-independence}

\begin{frame}{The independence assumption}
\protect\phantomsection\label{the-independence-assumption}
\begin{itemize}
\tightlist
\item
  What made the coin flip work? Treatment assignment was
  \textbf{independent of potential outcomes}: \[
  (Y^1, Y^0) \indep D
  \]
\item
  In words: knowing a bank's treatment status tells you \emph{nothing}
  about its potential outcomes --- neither its baseline \(Y^0\) nor its
  gain \(\delta\).
\item
  Under independence:

  \begin{itemize}
  \tightlist
  \item
    \(E[Y^0 \mid D{=}1] = E[Y^0 \mid D{=}0]\) \(\Rightarrow\)
    \textbf{selection bias \(= 0\)}
  \item
    \(E[\delta \mid D{=}1] = E[\delta \mid D{=}0]\) \(\Rightarrow\)
    \(ATT = ATU = ATE\) \(\Rightarrow\) \textbf{heterogeneity bias
    \(= 0\)}
  \end{itemize}
\item
  Both terms die, and \(SDO = ATE\). Randomization is the assignment
  mechanism that \emph{guarantees} independence.
\end{itemize}
\end{frame}

\begin{frame}{What randomization balances}
\protect\phantomsection\label{what-randomization-balances}
\begin{itemize}
\tightlist
\item
  Physical randomization balances treated and control groups on:

  \begin{itemize}
  \tightlist
  \item
    observable characteristics (size, leverage, ratings, \ldots),
  \item
    \textbf{unobservable} characteristics (management quality, risk
    appetite, \ldots),
  \item
    potential outcomes themselves, and
  \item
    treatment gains \(\delta_i\).
  \end{itemize}
\item
  No control variables needed. No model needed. The balance is created
  by the \emph{mechanism}, not by the econometrician.
\item
  Practice: in any experiment (or natural experiment), the first table
  is a \textbf{balance table} --- compare pre-treatment covariates
  across groups. Balance on observables is the testable \emph{symptom}
  of a sound mechanism.
\end{itemize}
\end{frame}

\begin{frame}{Case: eBay and the value of paid search}
\protect\phantomsection\label{case-ebay-and-the-value-of-paid-search}
\begin{itemize}
\tightlist
\item
  \textbf{The setting.} In paid search advertising, a firm bids so that
  its ad appears when a user searches a keyword. eBay spent heavily on
  this.
\item
  Two kinds of keyword:

  \begin{itemize}
  \tightlist
  \item
    \textbf{Brand} keywords contain the firm's own name --- ``eBay
    typewriter.''
  \item
    \textbf{Non-brand} keywords do not --- ``used typewriter.''
  \end{itemize}
\item
  \textbf{How the industry measured returns.} \emph{Attribution}: if a
  user clicks the ad and then buys, credit the sale to the ad. By this
  arithmetic, paid search looked extraordinarily profitable.
\item
  \textbf{The problem.} Somebody typing ``eBay typewriter'' was already
  going to eBay. Take the ad away and they click the organic eBay link
  one line below.
\item
  Attribution \textbf{conditions on clicking}. It never intervenes. It
  is an SDO, not an \(ATE\).
\end{itemize}
\end{frame}

\begin{frame}{eBay: what the experiments found}
\protect\phantomsection\label{ebay-what-the-experiments-found}
\begin{itemize}
\tightlist
\item
  Blake, Nosko, and Tadelis (2015, \emph{Econometrica}) ran large-scale
  field experiments at eBay, switching paid search off in randomly
  selected markets.
\item
  \textbf{Brand keywords}: \textbf{no measurable short-run benefit.}
  Traffic substituted almost entirely to the free organic link. eBay was
  paying for clicks it would have received anyway.
\item
  \textbf{Non-brand keywords}: ads \emph{did} work --- but only for
  \textbf{new and infrequent} users.

  \begin{itemize}
  \tightlist
  \item
    Frequent users were unaffected, yet accounted for most of the
    advertising spend.
  \item
    Net of costs, \textbf{average returns were negative}.
  \end{itemize}
\item
  Experimental returns were ``a fraction of conventional
  non-experimental estimates.''
\item
  \textbf{Two lessons.} Attribution measures selection, not causation.
  And the average effect concealed enormous heterogeneity --- precisely
  the variation in \(\delta_i\) we defined earlier.
\end{itemize}
\end{frame}

\begin{frame}{A genuine finance RCT: Karlan and Zinman (2009)}
\protect\phantomsection\label{a-genuine-finance-rct-karlan-and-zinman-2009}
\begin{itemize}
\tightlist
\item
  \emph{``Observing Unobservables: Identifying Information Asymmetries
  with a Consumer Credit Field Experiment,''} \textbf{Econometrica} 77.
\item
  \textbf{Setting.} A major South African consumer lender mailed loan
  offers to \textbf{58,000} former clients.
\item
  \textbf{Three independent randomizations:}

  \begin{enumerate}
  \tightlist
  \item
    The \textbf{offer rate} printed on the direct-mail solicitation.
  \item
    A \textbf{contract rate}, revealed only \emph{after} the client
    accepted the offer --- sometimes a surprise discount.
  \item
    A \textbf{dynamic repayment incentive}, also a surprise:
    preferential pricing on future loans for borrowers who stayed in
    good standing.
  \end{enumerate}
\item
  The surprises carry the design. Anything revealed \emph{after}
  acceptance cannot have influenced \textbf{who} accepted.
\end{itemize}
\end{frame}

\begin{frame}{Karlan and Zinman: separating two unobservables}
\protect\phantomsection\label{karlan-and-zinman-separating-two-unobservables}
\begin{itemize}
\tightlist
\item
  Credit markets have two classic private-information problems, and
  observational data cannot tell them apart:

  \begin{itemize}
  \tightlist
  \item
    \textbf{Adverse selection} (hidden \emph{type}): high rates attract
    intrinsically worse borrowers.
  \item
    \textbf{Moral hazard} (hidden \emph{action}): high rates weaken the
    incentive to repay.
  \end{itemize}
\item
  Both predict the same correlation --- higher rate, more default. No
  amount of loan data separates them.
\item
  The experiment does:

  \begin{itemize}
  \tightlist
  \item
    Vary the \textbf{offer} rate \(\Rightarrow\) changes \emph{who
    selects in} \(\Rightarrow\) identifies \textbf{selection}.
  \item
    Hold the offer rate fixed and vary the \textbf{contract} rate
    \(\Rightarrow\) same borrowers, different incentives \(\Rightarrow\)
    identifies \textbf{moral hazard}.
  \end{itemize}
\item
  \textbf{Findings}: strong evidence of moral hazard, weaker evidence of
  hidden information; roughly \textbf{13--21\%} of default attributable
  to moral hazard.
\item
  Randomization used not to estimate one number, but to
  \textbf{decompose a mechanism}.
\end{itemize}
\end{frame}

\begin{frame}{Other field experiments in finance}
\protect\phantomsection\label{other-field-experiments-in-finance}
\begin{itemize}
\tightlist
\item
  \textbf{Cole, Giné, and Vickery (2017, \emph{RFS})} --- rainfall
  insurance in India. Randomly insured farmers shifted into
  \textbf{higher-return, rainfall-sensitive cash crops}, \emph{before}
  rainfall was realized.
\item
  \textbf{Bertrand et al.~(2010, \emph{QJE})} --- direct-mail loan
  offers in South Africa, randomizing price \emph{and} content. A
  \textbf{photo of an attractive woman} raised demand as much as a
  \textbf{25\% rate cut}.
\item
  \textbf{Duflo and Saez (2003, \emph{QJE})} --- \$20 to attend a
  benefits fair raised enrollment, and raised it too for
  \textbf{untreated colleagues in treated departments}.
\item
  That last one is a clean randomization that still violates SUTVA (Part
  5).
\end{itemize}
\end{frame}

\begin{frame}{Where finance experiments live}
\protect\phantomsection\label{where-finance-experiments-live}
\begin{itemize}
\tightlist
\item
  Notice what these studies have in common: the researcher controls
  \textbf{an offer}.

  \begin{itemize}
  \tightlist
  \item
    A mailer, a price quote, an insurance policy, an invitation.
  \end{itemize}
\item
  That is why field experiments cluster in \textbf{household finance,
  credit, and insurance} --- markets with many small units and a lender
  or insurer willing to randomize.
\end{itemize}
\end{frame}

\begin{frame}{Why finance rarely gets to randomize}
\protect\phantomsection\label{why-finance-rarely-gets-to-randomize}
\begin{itemize}
\tightlist
\item
  Nobody will randomize leverage, M\&A, CEO pay, bankruptcy, or bank
  regulation:

  \begin{itemize}
  \tightlist
  \item
    stakes too high, units too big, ethics, general-equilibrium effects.
  \end{itemize}
\item
  So corporate finance and asset pricing live almost entirely in
  observational data.
\item
  \textbf{The rest of this course is a toolkit for \emph{earning}
  independence without a coin flip:}

  \begin{itemize}
  \tightlist
  \item
    Unconfoundedness / matching: independence \emph{conditional on
    observables}.
  \item
    RDD: independence \emph{near a threshold}.
  \item
    IV: independence \emph{of an instrument}.
  \item
    DiD / panel: independence \emph{of timing, in changes}.
  \end{itemize}
\item
  Every method = an argument that some comparison mimics the coin flip.
  Your job as a referee is to ask: \emph{whose coin, and who flipped
  it?}
\end{itemize}
\end{frame}

\section{SUTVA}\label{sutva}

\begin{frame}{The fine print under the potential outcomes}
\protect\phantomsection\label{the-fine-print-under-the-potential-outcomes}
\begin{itemize}
\tightlist
\item
  Writing \(Y_i^1, Y_i^0\) at all presumes the \textbf{Stable Unit
  Treatment Value Assumption}:
\end{itemize}

\begin{enumerate}
\tightlist
\item
  \textbf{No interference}: unit \(i\)'s potential outcomes do not
  depend on \emph{other units'} treatment status.
\item
  \textbf{No hidden versions of treatment}: ``treated'' means the same
  thing for every unit --- one treatment, one dose.
\end{enumerate}

\begin{itemize}
\tightlist
\item
  If SUTVA fails, \(Y_i^1\) is not even well-defined: bank \(i\)'s
  outcome under treatment depends on \emph{who else} was treated.
\end{itemize}
\end{frame}

\begin{frame}{Interference is the norm in finance}
\protect\phantomsection\label{interference-is-the-norm-in-finance}
\begin{itemize}
\tightlist
\item
  Units in finance \textbf{compete and interact}. Examples:

  \begin{itemize}
  \tightlist
  \item
    \textbf{Peer effects in capital structure} (Leary and Roberts 2014):
    firm \(i\)'s leverage responds to peers' leverage. Treating one firm
    treats its rivals.
  \item
    \textbf{Bailouts and competition}: a capital injection into bank
    \(i\) lets it bid for deposits and borrowers --- \emph{taking them
    from control banks}.
  \item
    \textbf{Banking DiDs}: treated banks' customers migrate to control
    banks; the control group's outcome moves \emph{because of} the
    treatment.
  \end{itemize}
\item
  If treatment hurts controls, the treated-control gap
  \textbf{overstates} the true effect; if it helps them (spillovers), it
  understates.
\end{itemize}
\end{frame}

\begin{frame}{``The control group is treated too''}
\protect\phantomsection\label{the-control-group-is-treated-too}
\begin{itemize}
\tightlist
\item
  In a competitive industry, there may be no untreated state of the
  world to observe:

  \begin{itemize}
  \tightlist
  \item
    Treatment effects estimated off treated-vs-control gaps are
    \textbf{relative} (partial-equilibrium) effects.
  \item
    The \textbf{aggregate} (general-equilibrium) effect can be smaller,
    zero, or of opposite sign.
  \end{itemize}
\item
  Example: a subsidy to treated firms raises their market share.
  Industry output may be unchanged --- pure reallocation. The DiD shows
  a large ``effect.''
\item
  Always ask: \emph{what happens to the market, not just the treated
  units?}
\item
  We will return to this in DiD (choosing clean controls) and in the
  reflection problem (peer effects).
\end{itemize}
\end{frame}

\begin{frame}{Hidden versions of treatment}
\protect\phantomsection\label{hidden-versions-of-treatment}
\begin{itemize}
\tightlist
\item
  SUTVA also requires the treatment to be \emph{one thing}.
\item
  Finance treatments are rarely uniform in dose:

  \begin{itemize}
  \tightlist
  \item
    ``Having a credit rating'' spans AAA to CCC.
  \item
    ``Being bailed out'' ranged from \$1bn to \$45bn in TARP.
  \item
    ``Adopting an anti-takeover provision'' bundles very different
    provisions.
  \end{itemize}
\item
  If doses differ and units select their dose, the single \(\delta_i\)
  hides a dose-response function --- and the estimand becomes an
  uninterpretable mixture.
\end{itemize}
\end{frame}

\begin{frame}{What to do about varying dose}
\protect\phantomsection\label{what-to-do-about-varying-dose}
\begin{itemize}
\tightlist
\item
  \textbf{Define} the treatment precisely. What exactly is the ``1'' in
  \(D_i = 1\)?
\item
  \textbf{Report} the \emph{distribution} of dose, not just a count of
  treated units.
\item
  If doses vary widely, do \textbf{not} collapse them into a binary.
  Instead, one of:

  \begin{itemize}
  \tightlist
  \item
    use the dose itself as the treatment and estimate a
    \textbf{dose-response function}, \(E[Y \mid \text{dose}]\);
  \item
    restrict the sample to a narrow dose band, so that ``treated'' means
    one thing;
  \item
    interact the treatment dummy with dose, and report the effect at
    several dose levels.
  \end{itemize}
\item
  Whichever you choose, state which estimand it delivers:
\end{itemize}

\begin{center}
\emph{``the effect of a \$1bn injection'' is not ``the effect of being bailed out.''}
\end{center}
\end{frame}

\section{The Regression Bridge}\label{the-regression-bridge}

\begin{frame}{Same problem, older dialect}
\protect\phantomsection\label{same-problem-older-dialect}
\begin{itemize}
\tightlist
\item
  You already know a language for causality --- the regression language:
  \[
  y = \beta_0 + \beta_1 x + u
  \]
\item
  Causal interpretation of \(\beta_1\) requires \textbf{conditional mean
  independence}: \[
  E[u \mid x] = E[u] = 0
  \]
\item
  The potential outcomes framework and the CMI framework describe the
  \textbf{same problem} in two dialects. Let's build the dictionary.
\end{itemize}
\end{frame}

\begin{frame}{The dictionary}
\protect\phantomsection\label{the-dictionary}
{\def\LTcaptype{none} % do not increment counter
\begin{longtable}[]{@{}
  >{\raggedright\arraybackslash}p{(\linewidth - 2\tabcolsep) * \real{0.5823}}
  >{\raggedright\arraybackslash}p{(\linewidth - 2\tabcolsep) * \real{0.4177}}@{}}
\toprule\noalign{}
\begin{minipage}[b]{\linewidth}\raggedright
Potential outcomes dialect
\end{minipage} & \begin{minipage}[b]{\linewidth}\raggedright
Regression dialect
\end{minipage} \\
\midrule\noalign{}
\endhead
\((Y^1, Y^0) \indep D\) & \(E[u \mid x] = E[u]\) (CMI) \\
Selection bias & \(\mathrm{cov}(x, u) \neq 0\) (OVB, simultaneity) \\
Heterogeneous \(\delta_i\) & Random coefficients \(b_i\) \\
ATE & Average partial effect (APE) \\
SDO & OLS coefficient of \(Y\) on \(D\) \\
\bottomrule\noalign{}
\end{longtable}
}

\begin{itemize}
\tightlist
\item
  The PO dialect is sharper about \emph{estimands} (ATE vs.~ATT) and
  \emph{assignment}.
\item
  The regression dialect is sharper about \emph{direction of bias} ---
  and it is the language referees speak. You need both.
\end{itemize}
\end{frame}

\begin{frame}{Omitted variable bias}
\protect\phantomsection\label{omitted-variable-bias}
\begin{itemize}
\tightlist
\item
  You estimate \(y = \beta_0 + \beta_1 x + u\), but the true model is \[
  y = \beta_0 + \beta_1 x + \beta_2 z + v
  \]
\item
  Then \[
  \hat{\beta}_1 \xrightarrow{p} \beta_1 + \underbrace{\frac{\mathrm{cov}(x, z)}{\mathrm{var}(x)}}_{\delta_{xz}} \beta_2
  \]
\item
  The bias is the \emph{effect of the omitted variable} (\(\beta_2\))
  times \emph{its regression on the included one} (\(\delta_{xz}\)).
\item
  Consistent only if \(\mathrm{cov}(x,z) = 0\) or \(\beta_2 = 0\): the
  omitted variable must matter \textbf{and} be correlated with \(x\).
\end{itemize}
\end{frame}

\begin{frame}{Signing the bias}
\protect\phantomsection\label{signing-the-bias}
\[
\text{bias} = \delta_{xz} \, \beta_2
\]

{\def\LTcaptype{none} % do not increment counter
\begin{longtable}[]{@{}lcc@{}}
\toprule\noalign{}
& \(\mathrm{cov}(x,z) > 0\) & \(\mathrm{cov}(x,z) < 0\) \\
\midrule\noalign{}
\endhead
\(\beta_2 > 0\) & \(+\) & \(-\) \\
\(\beta_2 < 0\) & \(-\) & \(+\) \\
\bottomrule\noalign{}
\end{longtable}
}

\begin{itemize}
\tightlist
\item
  Classic example: \(\ln(wage)\) on education, omitting ability.

  \begin{itemize}
  \tightlist
  \item
    Ability raises wages (\(\beta_2 > 0\)) and correlates positively
    with education \(\Rightarrow\) upward bias.
  \end{itemize}
\item
  Finance example: firm value on governance, omitting CEO talent.

  \begin{itemize}
  \tightlist
  \item
    Talented CEOs may sort into well-governed firms and raise value
    \(\Rightarrow\) governance coefficient too big.
  \end{itemize}
\item
  With \textbf{multiple} omitted variables (or multiple included
  regressors), signs interact and the algebra stops being your friend.
  Sign one bias at a time, or don't sign at all.
\end{itemize}
\end{frame}

\begin{frame}{The omitted variable is \(Y^0\)}
\protect\phantomsection\label{the-omitted-variable-is-y0}
\begin{itemize}
\tightlist
\item
  The switching equation factors as \(Y_i = Y_i^0 + D_i\delta_i\). Split
  each term into its mean plus a deviation, writing
  \(\delta = E[\delta_i] = ATE\): \[
  Y_i \;=\; \underbrace{E[Y^0]}_{\beta_0}
      \;+\; \underbrace{\delta}_{\beta_1} D_i
      \;+\; \underbrace{\left(Y_i^0 - E[Y^0]\right)}_{\text{baseline}}
      \;+\; \underbrace{D_i(\delta_i - \delta)}_{\text{gain}}
  \]
\item
  That \textbf{is} a regression of \(Y\) on a constant and \(D\) ---
  exactly, with no approximation and no assumptions. The last two terms
  are the error \(u_i\).
\item
  So the ``omitted variable'' is the \textbf{baseline potential outcome}
  \(Y_i^0\) (together with the individual gain \(\delta_i\)).
\item
  OVB --- in causal settings, it \emph{is} selection bias, written in
  the older dialect.
\item
  Careful: \(Y_i^0\) is \textbf{observed} for the untreated. It is
  counterfactual only for the treated --- missing exactly where we would
  need it. And no better data collection can fix that.
\end{itemize}
\end{frame}

\begin{frame}{Two ways treatment can be selected}
\protect\phantomsection\label{two-ways-treatment-can-be-selected}
\begin{itemize}
\tightlist
\item
  \(D\) can be correlated with \textbf{either} piece of the error: \[
  u_i = \underbrace{\left(Y_i^0 - E[Y^0]\right)}_{\text{baseline}} \;+\; \underbrace{D_i(\delta_i - \delta)}_{\text{gain}}
  \]
\item
  \textbf{Selection on levels} --- treatment goes to units with unusual
  \(Y_i^0\): how they would have done \emph{anyway}, untreated.

  \begin{itemize}
  \tightlist
  \item
    Bailing out the weakest banks. Training the least employable
    workers.
  \end{itemize}
\item
  \textbf{Selection on gains} --- treatment goes to units with unusual
  \(\delta_i\): how much they would \emph{benefit}.

  \begin{itemize}
  \tightlist
  \item
    The Perfect Regulator. Any optimizing agent allocating a scarce
    treatment.
  \end{itemize}
\item
  The two are distinct, and \textbf{either one alone} breaks the naive
  comparison.
\item
  Note our own running example sorts on \textbf{gains}, not levels.
  Defining selection bias as ``\(\mathrm{cov}(D,Y^0) \neq 0\)'' would
  misclassify it.
\end{itemize}
\end{frame}

\begin{frame}{Both biases, from one lemma}
\protect\phantomsection\label{both-biases-from-one-lemma}
\begin{itemize}
\tightlist
\item
  \textbf{Lemma.} For \emph{binary} \(D\) with \(\pi = P(D{=}1)\), and
  any random variable \(W\): \[
  \mathrm{cov}(D, W) = \pi(1-\pi)\Big(E[W \mid D{=}1] - E[W \mid D{=}0]\Big)
  \]

  \begin{itemize}
  \tightlist
  \item
    \emph{Proof.} \(E[DW] = \pi E[W|D{=}1]\) since \(D\) is binary, and
    \(E[W] = \pi E[W|D{=}1] + (1-\pi)E[W|D{=}0]\). Substitute into
    \(\mathrm{cov}(D,W) = E[DW] - E[D]E[W]\) and collect.
    \(\blacksquare\)
  \end{itemize}
\item
  \textbf{Take \(W = Y^0\):} \[
  \mathrm{cov}(D, Y^0) = \pi(1-\pi)\Big(E[Y^0|D{=}1] - E[Y^0|D{=}0]\Big)
  \]
\item
  \textbf{Take \(W = D_i(\delta_i - \delta)\).} It is \(0\) whenever
  \(D=0\), so \(E[W|D{=}0] = 0\), while \(E[W|D{=}1] = ATT - ATE\): \[
  \mathrm{cov}\big(D,\, D_i(\delta_i-\delta)\big) = \pi(1-\pi)\big(ATT - ATE\big)
  \]
\item
  Since \(\mathrm{var}(D) = \pi(1-\pi)\), dividing each by
  \(\mathrm{var}(D)\) \textbf{cancels that factor} --- leaving the
  selection bias term and \(ATT - ATE\).
\end{itemize}
\end{frame}

\begin{frame}{The two dialects give the same equation}
\protect\phantomsection\label{the-two-dialects-give-the-same-equation}
\begin{itemize}
\tightlist
\item
  Since
  \(\hat\beta_1 \xrightarrow{p} \delta + \mathrm{cov}(D,u)/\mathrm{var}(D)\),
  and using \(ATT - ATE = (1-\pi)(ATT-ATU)\) from earlier: \[
  \hat\beta_1 \;\xrightarrow{p}\;
  \underbrace{ATE}_{\delta}
  \;+\; \underbrace{\Big(E[Y^0|D{=}1]-E[Y^0|D{=}0]\Big)}_{\text{selection bias}}
  \;+\; \underbrace{(1-\pi)(ATT-ATU)}_{\text{heterogeneity bias}}
  \]
\item
  The right-hand side is exactly the \(SDO\) decomposition we proved in
  Part 3.
\item
  The regression route and the potential-outcomes route are not
  analogies. \textbf{They are the same three terms.}
\item
  So \(\hat\beta_1\) is inconsistent for the \(ATE\) --- but perfectly
  \emph{consistent} for the \(SDO\). Once again: OLS delivers the
  estimand you asked for, not the one you wanted.
\end{itemize}
\end{frame}

\begin{frame}{Simultaneity: when \(y\) fires back}
\protect\phantomsection\label{simultaneity-when-y-fires-back}
\begin{itemize}
\tightlist
\item
  Suppose the regressor responds to the outcome: \[
  y = \beta_0 + \beta_1 x + u, \qquad x = \alpha_0 + \alpha_1 y + \varepsilon
  \]
\item
  Solve for \(x\): \[
  x = \frac{\alpha_0 + \alpha_1 \beta_0}{1 - \alpha_1 \beta_1} + \frac{\alpha_1 u + \varepsilon}{1 - \alpha_1 \beta_1}
  \]
\item
  \(x\) is mechanically correlated with \(u\): CMI fails, OLS is biased
  --- \textbf{reverse causality}.
\item
  Finance examples: leverage and performance; board independence and
  performance; analyst coverage and returns. In equilibrium, almost
  every firm characteristic responds to almost every outcome.
\item
  Note: lagging \(x\) does \textbf{not} fix this if \(x\) is persistent
  and anticipates \(y\). It hides the problem; it does not solve it.
\end{itemize}
\end{frame}

\begin{frame}{Random coefficients: the formal bridge}
\protect\phantomsection\label{random-coefficients-the-formal-bridge}
\begin{itemize}
\tightlist
\item
  Let every firm have its own line: \[
  y_i = a_i + b_i x_i, \qquad a_i = \alpha + c_i, \quad b_i = \beta + d_i
  \]
\item
  With binary \(x\), this \textbf{is} the potential outcomes model:
  \(b_i = \delta_i\), \(a_i = Y_i^0\).
\item
  Rewrite:
  \(y_i = \alpha + \beta x_i + \underbrace{(c_i + d_i x_i)}_{u_i}\)
\item
  OLS identifies the \textbf{average partial effect} \(\beta = E[b_i]\)
  (the ATE!) only if \[
  E[c_i \mid x_i] = 0 \;\;(\text{no selection on levels}), \qquad E[d_i \mid x_i] = 0 \;\;(\text{no selection on gains})
  \]
\item
  These are precisely ``selection bias \(=0\)'' and ``ATT \(=\) ATU.''
  The two dialects have converged.
\end{itemize}
\end{frame}

\begin{frame}{Bad controls: a preview}
\protect\phantomsection\label{bad-controls-a-preview}
\begin{itemize}
\tightlist
\item
  Adding controls is not free. If a control \(x_2\) is itself
  \textbf{affected by the treatment}: \[
  y = \beta_0 + \beta_1 x_1 + \beta_2 x_2 + u, \qquad x_2 = \gamma_0 + \gamma_1 x_1 + \varepsilon
  \]
\item
  then conditioning on \(x_2\) blocks part of the effect you want ---
  and can \emph{introduce} selection bias where none existed.
\item
  Rule of thumb: controls must be \textbf{pre-treatment} (determined
  before \(D\)).
\item
  Why conditioning on a post-treatment variable creates bias is subtle
  --- it is the centerpiece of the next lecture (DAGs and colliders).
\end{itemize}
\end{frame}

\section{Randomization Inference}\label{randomization-inference}

\begin{frame}{Two null hypotheses}
\protect\phantomsection\label{two-null-hypotheses}
\begin{itemize}
\tightlist
\item
  Standard (Neyman) null: the \textbf{average} effect is zero,
  \(E[\delta_i] = 0\).

  \begin{itemize}
  \tightlist
  \item
    Inference relies on asymptotics: large \(n\), estimated standard
    errors.
  \end{itemize}
\item
  Fisher's \textbf{sharp null}: the effect is zero \emph{for every
  unit}, \[
  H_0: \delta_i = 0 \;\; \text{for all } i
  \]
\item
  Under the sharp null, both potential outcomes are known for every
  unit: \(Y_i^1 = Y_i^0 = Y_i\) (the observed value).
\item
  That makes \emph{exact} inference possible --- no asymptotics, no
  standard errors, any sample size.
\end{itemize}
\end{frame}

\begin{frame}{The permutation logic}
\protect\phantomsection\label{the-permutation-logic}
\begin{enumerate}
\tightlist
\item
  Under the sharp null, the observed outcomes would have been the same
  under \textbf{any} treatment assignment.
\item
  So: re-randomize \(D\) on paper, in every way the actual randomization
  could have come out.
\item
  Recompute the test statistic (e.g., the SDO) for each hypothetical
  assignment.
\item
  The exact p-value is the share of assignments giving a statistic at
  least as extreme as the observed one.
\end{enumerate}

\begin{itemize}
\tightlist
\item
  The randomness comes from the \textbf{assignment mechanism itself} ---
  the same mechanism that justified the causal comparison justifies the
  inference.
\end{itemize}
\end{frame}

\begin{frame}[fragile]{Randomization inference on ten banks}
\protect\phantomsection\label{randomization-inference-on-ten-banks}
\begin{itemize}
\tightlist
\item
  Suppose the ten banks were randomized: five drawn by lot got
  injections. Observed SDO \(= 1.6\).
\end{itemize}

\begin{Shaded}
\begin{Highlighting}[]
\NormalTok{y }\OtherTok{\textless{}{-}} \FunctionTok{c}\NormalTok{(}\DecValTok{2}\NormalTok{, }\DecValTok{6}\NormalTok{, }\DecValTok{7}\NormalTok{, }\DecValTok{2}\NormalTok{, }\DecValTok{6}\NormalTok{, }\DecValTok{6}\NormalTok{, }\DecValTok{7}\NormalTok{, }\DecValTok{7}\NormalTok{, }\DecValTok{7}\NormalTok{, }\DecValTok{8}\NormalTok{)     }\CommentTok{\# observed lending growth}
\NormalTok{treated }\OtherTok{\textless{}{-}} \FunctionTok{c}\NormalTok{(}\DecValTok{2}\NormalTok{, }\DecValTok{3}\NormalTok{, }\DecValTok{5}\NormalTok{, }\DecValTok{8}\NormalTok{, }\DecValTok{9}\NormalTok{)              }\CommentTok{\# banks drawn by lot}
\NormalTok{sdo\_obs }\OtherTok{\textless{}{-}} \FunctionTok{mean}\NormalTok{(y[treated]) }\SpecialCharTok{{-}} \FunctionTok{mean}\NormalTok{(y[}\SpecialCharTok{{-}}\NormalTok{treated])}

\NormalTok{perms }\OtherTok{\textless{}{-}} \FunctionTok{combn}\NormalTok{(}\DecValTok{10}\NormalTok{, }\DecValTok{5}\NormalTok{)                     }\CommentTok{\# all 252 possible assignments}
\NormalTok{sdo\_all }\OtherTok{\textless{}{-}} \FunctionTok{apply}\NormalTok{(perms, }\DecValTok{2}\NormalTok{, }\ControlFlowTok{function}\NormalTok{(tr) }\FunctionTok{mean}\NormalTok{(y[tr]) }\SpecialCharTok{{-}} \FunctionTok{mean}\NormalTok{(y[}\SpecialCharTok{{-}}\NormalTok{tr]))}
\NormalTok{p\_exact }\OtherTok{\textless{}{-}} \FunctionTok{mean}\NormalTok{(}\FunctionTok{abs}\NormalTok{(sdo\_all) }\SpecialCharTok{\textgreater{}=} \FunctionTok{abs}\NormalTok{(sdo\_obs))}
\FunctionTok{round}\NormalTok{(}\FunctionTok{c}\NormalTok{(}\AttributeTok{SDO =}\NormalTok{ sdo\_obs, }\AttributeTok{exact\_p =}\NormalTok{ p\_exact), }\DecValTok{3}\NormalTok{)}
\end{Highlighting}
\end{Shaded}

\begin{verbatim}
    SDO exact_p 
  1.600   0.397 
\end{verbatim}

\begin{itemize}
\tightlist
\item
  With \(n = 10\), a difference of 1.6pp is \emph{not} statistically
  distinguishable from luck-of-the-draw --- and we know this
  \textbf{exactly}, not asymptotically.
\end{itemize}
\end{frame}

\begin{frame}{Where you will see this again}
\protect\phantomsection\label{where-you-will-see-this-again}
\begin{itemize}
\tightlist
\item
  Randomization inference is not a curiosity:

  \begin{itemize}
  \tightlist
  \item
    \textbf{Synthetic control} (week 11): placebo tests permute the
    treatment across donor units --- that \emph{is} randomization
    inference.
  \item
    \textbf{Small-N settings}: one treated state, one regulatory event,
    a handful of treated firms --- common in finance, hopeless for
    asymptotics.
  \item
    Cluster-robust inference with few clusters: permutation tests as the
    honest fallback.
  \end{itemize}
\item
  The deeper point: \emph{inference should mirror the assignment
  mechanism.} Where the mechanism is uncertain, so is the p-value.
\end{itemize}
\end{frame}

\section{Wrapping Up}\label{wrapping-up}

\begin{frame}{Summary {[}Part 1{]}}
\protect\phantomsection\label{summary-part-1}
\begin{itemize}
\tightlist
\item
  Causality is defined by \textbf{counterfactuals}; the fundamental
  problem is that one potential outcome is always missing.
\item
  Estimands differ: \textbf{ATE, ATT, ATU} answer different policy
  questions. Know which one your design identifies.
\item
  The naive comparison decomposes as \[
  SDO = ATE + \text{selection bias} + (1-\pi)(ATT - ATU)
  \]
\item
  Optimizing agents generate selection \textbf{on gains} --- the Perfect
  Regulator gets the sign wrong \emph{because} it allocates well.
\item
  More data does not fix a biased assignment mechanism.
\end{itemize}
\end{frame}

\begin{frame}{Summary {[}Part 2{]}}
\protect\phantomsection\label{summary-part-2}
\begin{itemize}
\tightlist
\item
  \textbf{Randomization} kills both bias terms by making
  \((Y^1, Y^0) \indep D\). Balance tables are its testable symptom.
\item
  \textbf{SUTVA} fails easily in finance: competition, peer effects, and
  heterogeneous doses. Ask what happens to the market, not just the
  treated.
\item
  The \textbf{regression dialect} translates one-for-one: CMI
  \(\Leftrightarrow\) independence; OVB \(\Leftrightarrow\) selection
  bias; random coefficients \(\Leftrightarrow\) heterogeneous
  \(\delta_i\).
\item
  \textbf{Randomization inference} gives exact p-values from the
  assignment mechanism --- remember it when you meet synthetic control.
\item
  Next lecture: DAGs --- a graphical language for deciding \emph{what to
  control for, and what not to}.
\end{itemize}
\end{frame}

\begin{frame}{References}
\protect\phantomsection\label{references}
\footnotesize

\begin{itemize}
\tightlist
\item
  Cunningham, S. \emph{Causal Inference: The Remix}, Ch. 4.
\item
  Holland, P. (1986). Statistics and Causal Inference. \emph{JASA} 81,
  945--960.
\item
  Blake, T., C. Nosko, and S. Tadelis (2015). Consumer Heterogeneity and
  Paid Search Effectiveness: A Large-Scale Field Experiment.
  \emph{Econometrica} 83, 155--174.
\item
  Karlan, D. and J. Zinman (2009). Observing Unobservables: Identifying
  Information Asymmetries with a Consumer Credit Field Experiment.
  \emph{Econometrica} 77, 1993--2008.
\item
  Bertrand, M., D. Karlan, S. Mullainathan, E. Shafir, and J. Zinman
  (2010). What's Advertising Content Worth? \emph{Quarterly Journal of
  Economics} 125, 263--306.
\item
  Cole, S., X. Giné, and J. Vickery (2017). How Does Risk Management
  Influence Production Decisions? \emph{Review of Financial Studies} 30,
  1935--1970.
\item
  Duflo, E. and E. Saez (2003). The Role of Information and Social
  Interactions in Retirement Plan Decisions. \emph{Quarterly Journal of
  Economics} 118, 815--842.
\item
  Leary, M. and M. Roberts (2014). Do Peer Firms Affect Corporate
  Financial Policy? \emph{Journal of Finance} 69, 139--178.
\item
  Angrist, J. and J.-S. Pischke (2009). \emph{Mostly Harmless
  Econometrics}. Princeton UP.
\end{itemize}
\end{frame}




\end{document}
